% ==============================================================
% ARO-KONFORMES REVIEW -- TEIL 1
% Projektive Strukturtheorie (PST) -- Kritische und Bedeutsame Befunde
% Reviewer: Kimi (Moonshot AI)
% Datum: Juli 2026
% Version: 1.1 (korrigiert und aufgeteilt)
% ==============================================================

\documentclass[11pt,a4paper]{article}

\usepackage[utf8]{inputenc}
\usepackage[T1]{fontenc}
\usepackage[ngerman]{babel}
\usepackage{amsmath,amssymb,amsthm}
\usepackage{geometry}
\usepackage{parskip}
\usepackage{enumitem}
\usepackage{booktabs}
\usepackage{xcolor}
\usepackage{hyperref}
\usepackage{longtable}

\geometry{margin=2.5cm}
\setlist{itemsep=0.3em, topsep=0.3em}
\hypersetup{
  colorlinks=true,
  linkcolor=blue,
  urlcolor=blue,
  pdftitle={ARO-Review: PST Kritische und Bedeutsame Befunde},
  pdfauthor={Kimi -- Forschungsprogramm Ontophänomenalismus}
}

\newcommand{\K}[1]{\textbf{[K-#1]}}
\newcommand{\B}[1]{\textbf{[B-#1]}}
\newcommand{\M}[1]{\textbf{[M-#1]}}
\newcommand{\Pos}[1]{\textbf{[P-#1]}}

\newcommand{\Dok}[1]{\noindent\textbf{Dokument:} #1\par\smallskip}
\newcommand{\Stelle}[1]{\noindent\textbf{Stelle:} #1\par\smallskip}
\newcommand{\Prob}[1]{\noindent\textbf{Problem:} #1\par\smallskip}
\newcommand{\Krit}[1]{\noindent\textbf{Warum kritisch:} #1\par\smallskip}
\newcommand{\Empf}[1]{\noindent\textbf{Empfehlung:} #1\par\smallskip}

\newenvironment{Fix}
  {\par\medskip\noindent\textbf{Fixing-Vorschlag:}\par\nobreak\smallskip}
  {\par\medskip}

\title{
  \textbf{Adversarial Review: Projektive Strukturtheorie (PST)}\\[0.4em]
  {\large Teil~1: Kritische und Bedeutsame Befunde}\\[0.5em]
  {\normalsize ARO-Konform -- Version~1.1}
}
\author{
  \textbf{Reviewer: Kimi}\\
  \small Forschungsprogramm Ontophänomenalismus
}
\date{Juli 2026}

\begin{document}

\maketitle

\begin{abstract}
Dieses Review bewertet den vollständigen PST-Dokumentenzyklus
(PST-1-F bis PST-6-R) gemäß dem Adversarial Review Protocol (AR).
Dieser erste Teil enthält die \textbf{kritischen Befunde [K]} und
\textbf{bedeutsamen Befunde [B]}. Alle Befunde sind mit konkreten
Fixing-Vorschlägen versehen. Die \textbf{minoren Befunde [M]},
\textbf{positiven Bewertungen [P]} sowie die \textbf{priorisierten
Empfehlungen} finden sich in Teil~2.
\end{abstract}

\tableofcontents
\newpage

%==============================================================
\section{Zusammenfassung des Gesamteindrucks}
%==============================================================

Die PST ist ein \textbf{methodisch diszipliniertes, philosophisch
kohärentes Forschungsprogramm} im Sinne von Lakatos. Der harte Kern
(OPP als prä-metrischer Urgrund, Projektion als Erzeugung von
Raumzeit, minimale Ontologie) ist konsistent und anschlussfähig an
die philosophischen Module des OP.

Die größten Stärken liegen in der \textbf{methodischen
Ehrlichkeit} (No-Go-Ergebnisse werden als Ressource genutzt), der
\textbf{Selbstkritik} (Rangneutralität, Zirkularitätstests) und
der \textbf{klaren Priorisierung} (PST-5-S).

Die größten Schwächen liegen in drei Bereichen:
\begin{enumerate}
  \item \textbf{Mathematische Inkonsistenz:} Der gewählte Weg zur
    Lorentz-Signatur (komplexe Gram-Matrizen, Wick-Rotation,
    imaginäre Eigenwerte) ist mit der in PST-2-M festgelegten
    reell-symmetrischen Gram-Matrix unvereinbar.
  \item \textbf{Empirische Undokumentiertheit:} Die zentrale
    Behauptung $d_{\text{eff}} \approx 4$ bei $N = 10\,000$ ist
    nicht reproduzierbar, da Code und Rohdaten fehlen.
  \item \textbf{Fehlende philosophische Stärke:} Die originellste
    Idee des gesamten OP-Programms -- die konstruktiven
    Raumdimensionen vs. die destruktive Zeitdimension als
    ontologische Begründung der Metrik-Signatur -- taucht in
    keinem PST-Dokument auf.
\end{enumerate}

\newpage

%==============================================================
\section{Kritische Befunde [K]}
%==============================================================

%--------------------------------------------------------------
\subsection{\K{01} PST-1-F: ``Existenz = Bestimmtheit'' ist ein petitio principii}
%--------------------------------------------------------------

\Dok{PST-1-F, Sektion 3}
\Stelle{Theorem ``Existenz = Bestimmtheit''}
\Prob{Der Beweis ist ein Widerspruchsbeweis: Angenommen, X existiert,
aber ist vollkommen unbestimmt. Dann ist X nicht vom Nichts
unterscheidbar. Aber ``existiert'' in der Annahme setzt bereits
Bestimmtheit voraus (per UP).}
\Krit{Dies ist kein Widerspruchsbeweis im strengen Sinne, sondern eine
\textbf{petitio principii}: Die Konklusion ist in die Prämisse
eingebaut. Wenn man die UP ablehnt, bricht der Beweis zusammen. Das
Theorem suggeriert einen höheren Formalitätsgrad als
vorhanden.}
\Empf{Umbenennen in ``Proposition'' oder ``Arbeitshypothese''. Der
Beweis ist gültig \emph{innerhalb} des OP-Systems, aber nicht
absolut.}
\begin{Fix}
In PST-1-F, Sektion 3, die Theorem-Umgebung durch eine
\texttt{hypothese}-Umgebung ersetzen und hinzufügen: ``Diese
Proposition ist gültig innerhalb des OP-Axiomensystems (UP als
Axiom). Sie ist kein absoluter Beweis der Gleichsetzung von Existenz
und Bestimmtheit.''
\end{Fix}

%--------------------------------------------------------------
\subsection{\K{02} PST-2-M: ``Theorem'' zur Planck-Schranke ist kein Theorem}
%--------------------------------------------------------------

\Dok{PST-2-M, Sektion 7}
\Stelle{Theorem ``Planck-Schranke für bestimmte Existenz''}
\Prob{Das Dokument präsentiert eine Behauptung als
\texttt{theorem}-Umgebung:
\[
  \mathcal{B}(X,Y) = 1 \iff MI(X,Y) \geq I_0 \iff d_{XY} \geq \ell_P
\]
Die Methodennote sagt im selben Atemzug: ``Dieses Theorem ist eine
Arbeitshypothese ... noch nicht mathematisch bewiesen.''}
\Krit{In LaTeX (und in der Mathematik) ist eine \texttt{theorem}-Umgebung
eine \textbf{Behauptung mit Beweis}. Eine Arbeitshypothese gehört in
eine \texttt{hypothese}- oder \texttt{conjecture}-Umgebung. Die Mischung
von \texttt{theorem} mit ``ist eine Arbeitshypothese'' ist ein
\textbf{Selbstwiderspruch auf der Präsentationsebene}. Ein externer
Mathematiker würde hier sofort aufhören zu lesen.}
\Empf{In zwei Schritte aufteilen: (1) \texttt{lemma} oder
\texttt{proposition}: $MI(X,Y) \geq I_0 \iff \mathcal{B}(X,Y) = 1$
(Definition/Operationalisierung). (2) \texttt{hypothese}: Diese
Bestimmtheit korrespondiert mit $d_{XY} \geq \ell_P$
(physikalische Brücke).}
\begin{Fix}
In PST-2-M, Sektion 7, die \texttt{theorem}-Umgebung streichen und
durch zwei separate Umgebungen ersetzen: Eine
\texttt{definition}/\texttt{proposition} für die
Informationsbedingung und eine \texttt{hypothese} für die
Planck-Korrespondenz.
\end{Fix}

%--------------------------------------------------------------
\subsection{\K{03} PST-2-M: Die Verbindung zwischen \textbf{R} und \textbf{K} fehlt}
%--------------------------------------------------------------

\Dok{PST-2-M, Sektion 3 und 4}
\Stelle{Definition von \textbf{R} und Definition von \textbf{K}}
\Prob{Es werden zwei Matrizen definiert:
\begin{itemize}
  \item $\mathbf{R} \in \mathbb{R}^{N \times N}$ mit $R_{ij} = r(s_i, s_j)$,
    wobei $r: S \times S \to \mathbb{R}_{\geq 0}$ eine ``Relationsstärke'' ist.
  \item $\mathbf{K} \in \mathbb{R}^{N \times N}$ mit $K_{ij} = k(s_i, s_j)$,
    wobei $k$ ein ``symmetrisch-positiv semidefiniter Kernel'' ist.
\end{itemize}
Aber: \textbf{Wie hängen $\mathbf{R}$ und $\mathbf{K}$ zusammen?}}
\Krit{Die Hypothese in Sektion 3 suggeriert: $r(s_i, s_j) = MI(X_i; X_j)$.
Aber in Sektion 4 wird $K_{ij} = k(s_i, s_j)$ definiert, ohne Bezug zu
$\mathbf{R}$. Wenn $\mathbf{K} \neq \mathbf{R}$, dann ist die gesamte
Kette ``Relation $\to$ Kernel $\to$ Spektrum $\to$ Geometrie''
unterbrochen. Wenn $\mathbf{K} = \mathbf{R}$, dann muss $\mathbf{R}$
symmetrisch-positiv semidefinit sein -- was bei einer allgemeinen
Relationsstärke $r: S \times S \to \mathbb{R}_{\geq 0}$ \textbf{nicht
gegeben} ist.}
\Empf{Explizit definieren: Entweder $\mathbf{K} = \mathbf{R}$ und
$\mathbf{R}$ muss als Kernel konstruiert werden, oder $\mathbf{R}$ ist
die Rohdaten-Matrix und $\mathbf{K} = \varphi(\mathbf{R})$ ist eine
Kernel-Konstruktion.}
\begin{Fix}
In PST-2-M, Sektion 4, vor der Definition von $\mathbf{K}$ eine
Zwischendefinition einfügen: ``Die Gram-Matrix $\mathbf{K}$ wird aus
der Relationsmatrix $\mathbf{R}$ durch einen positiv-semidefiniten
Kernel konstruiert: $K_{ij} = \exp(-R_{ij}^2 / 2\sigma^2)$
(Gaußscher Kernel) oder $K_{ij} = R_{ij}$ falls $\mathbf{R}$ bereits
positiv semidefinit ist.''
\end{Fix}

%--------------------------------------------------------------
\subsection{\K{04} PST-2-M: ``Isometrische Einbettung'' ist mathematisch irreführend}
%--------------------------------------------------------------

\Dok{PST-2-M, Sektion 4}
\Stelle{Theorem ``Spektralzerlegung der Gram-Matrix''}
\Prob{Das Theorem behauptet:
\[
  \Phi: s_i \mapsto (\sqrt{\lambda_1} v_{1i}, \ldots, \sqrt{\lambda_N} v_{Ni})
\]
sei ``eine isometrische Einbettung der Zustände in den euklidischen
Raum $\mathbb{R}^N$''.}
\Krit{Eine isometrische Einbettung erhält Abstände. Aber:
\begin{itemize}
  \item Der OPP hat \textbf{keine Metrik} (per Definition).
  \item Die Einbettung $\Phi$ ist die Kernel-PCA-Einbettung. Sie erhält
    das Skalarprodukt im Feature-Raum (definiert durch $\mathbf{K}$),
    aber nicht einen ``Abstand'' im OPP -- denn einen solchen gibt es
    nicht.
\end{itemize}
Das Wort ``isometrisch'' suggeriert, dass die Einbettung eine
vorgegebene Metrik erhält. Aber genau das lehnt die PST ab. Es ist
ein \textbf{Kategorienfehler}: Die PST behauptet, Metrik entstehe erst
durch die Einbettung -- aber dann kann die Einbettung nicht
``isometrisch'' sein.}
\Empf{Formulieren als: ``Die Einbettung $\Phi$ induziert ein
Skalarprodukt und damit eine Metrik auf dem Bild der Zustände in
$\mathbb{R}^N$. Diese Metrik ist emergent -- sie existiert nicht im
OPP, sondern erst in der projizierten Darstellung.''}
\begin{Fix}
In PST-2-M, Sektion 4, den Satz ``eine isometrische Einbettung
...'' ersetzen durch: ``Die Einbettung $\Phi$ repräsentiert die
durch $\mathbf{K}$ definierte Ähnlichkeitsstruktur im euklidischen
Raum $\mathbb{R}^N$ und induziert dort eine emergente Metrik.''
\end{Fix}

%--------------------------------------------------------------
\subsection{\K{05} PST-3-P: Die Stabilitätshypothese für $d = 4$ ist zirkulär}
%--------------------------------------------------------------

\Dok{PST-3-P, Sektion 4}
\Stelle{Hypothese ``Stabilitätshypothese für $d = 4$''}
\Prob{Die Hypothese begründet $d = 4$ mit physikalischen Argumenten
aus der \emph{bekannten} Physik:
\begin{itemize}
  \item Planetenbahnen sind nur in $d = 3$ stabil
  \item Yang-Mills-Theorien sind in $d = 4$ renormierbar
  \item $SO(1,3)$ hat in $d = 4$ Weyl-Spinoren
\end{itemize}}
\Krit{Die PST soll die Physik \emph{aus dem OPP herleiten}, nicht die
bekannte Physik als Begründung für $d = 4$ verwenden. Wenn ihr
sagt: ``$d = 4$ emergiert, weil in $d = 4$ Planetenbahnen stabil
sind'', dann setzt ihr voraus, was ihr beweisen wollt. Das ist ein
\textbf{petitio principii} auf der Ebene der Physik.}
\Empf{Entweder (a) die physikalischen Argumente als \emph{heuristische
Randbedingungen} kennzeichnen (``Wir wissen aus der bekannten Physik,
dass $d = 4$ funktioniert; die PST muss zeigen, \emph{warum} der OPP
notwendigerweise $d = 4$ liefert'') -- oder (b) die Argumente durch
\emph{strukturelle} Argumente ersetzen, die unabhängig von der
bekannten Physik sind (z.B. aus der Spektraltheorie der Gram-Matrix,
aus der Stabilität von Diffusion Maps, aus der
Informationsgeometrie).}
\begin{Fix}
In PST-3-P, Sektion 4, die Hypothese umbenennen in
``Heuristische Randbedingung'' und formulieren: ``Die bekannte Physik
liefert die Randbedingung $d = 4$ für physikalische
Projektionsoperatoren. Die PST sucht einen Operator, der diese
Randbedingung erfüllt, ohne sie vorauszusetzen.''
\end{Fix}

%--------------------------------------------------------------
\subsection{\K{06} PST-3-P: Die ``Definition'' von $\mathcal{P}^*$ ist eine Wunschliste}
%--------------------------------------------------------------

\Dok{PST-3-P, Sektion 5}
\Stelle{Definition ``Universeller Projektionsoperator''}
\Prob{Ein mathematisches Objekt wird definiert durch seine
\emph{Eigenschaften}, nicht durch seine \emph{Ziele}. Die
``Definition'' listet fünf Bedingungen auf, die $\mathcal{P}^*$
erfüllen soll -- aber sie sagt nicht, \emph{was} $\mathcal{P}^*$
ist. Ist es eine Abbildung? Ein Funktor? Ein Operator auf einem
Hilbert-Raum? Ein stochastischer Kern?}
\Krit{Ohne mathematische Spezifikation ist $\mathcal{P}^*$ ein
Platzhalter, kein Objekt. Das ist für ein Forschungsprogramm in
Ordnung -- aber es darf nicht als \texttt{definition} in einem
mathematischen Dokument erscheinen.}
\Empf{Umbenennen in \texttt{hypothese} oder \texttt{desideratum}. Die
fünf Bedingungen sind \emph{Randbedingungen}, keine Definition.}
\begin{Fix}
In PST-3-P, Sektion 5, die \texttt{definition}-Umgebung durch
\texttt{hypothese} oder ein neues \texttt{desideratum} ersetzen. Die
fünf Bedingungen als ``Randbedingungen für $\mathcal{P}^*$''
kennzeichnen. Hinzufügen: ``Eine vollständige mathematische
Definition von $\mathcal{P}^*$ erfordert die in PST-5-S, Priorität 3
geforderte Fundierung (Kategorientheorie, Operatoralgebren oder
spektrale Tripel).''
\end{Fix}

%--------------------------------------------------------------
\subsection{\K{07} PST-3-P: Lorentz-Signatur aus ``Imaginärteil von Eigenvektoren'' ist mathematisch problematisch}
%--------------------------------------------------------------

\Dok{PST-3-P, Sektion 7}
\Stelle{Hypothese ``Lorentz-Signatur aus Projektion''}
\Prob{Die Hypothese suggeriert, dass die Minkowski-Signatur
$(+,-,-,-)$ aus dem Imaginärteil von Eigenvektoren bei komplexen
Kernel-Matrizen emergieren könnte. Aber:
\begin{itemize}
  \item In PST-2-M ist die Gram-Matrix \textbf{reell-symmetrisch}
    definiert: $K_{ij} = k(s_i, s_j) \in \mathbb{R}$.
  \item Reell-symmetrische Matrizen haben \textbf{reelle Eigenwerte und
    reelle Eigenvektoren}. Es gibt keinen Imaginärteil.
  \item Wenn man zu komplexen Kerneln übergeht, verliert man die
    Symmetrie $\mathbf{K} = \mathbf{K}^T$ und damit die Garantie
    reeller Eigenwerte. Dann ist die Spektralzerlegung nicht mehr
    garantiert reell.
\end{itemize}}
\Krit{Das ist ein \textbf{mathematischer Kategorienfehler}. Man kann
nicht aus einer reellen Struktur einen Imaginärteil zaubern, der
dann die Lorentz-Signatur erzeugt. Das wäre, als wollte man aus einem
reellen Vektorraum eine komplexe Struktur extrahieren -- das geht
nur, wenn man eine komplexe Struktur \emph{voraussetzt}.}
\Empf{Entweder (a) die Gram-Matrix von vorneherein als hermitesch
(komplex) definieren -- was die gesamte PST-2-M ändern würde --
oder (b) einen anderen Mechanismus für die Lorentz-Signatur suchen.
Ein vielversprechender Ansatz wäre: Die Signatur emergiert aus der
\textbf{Ordnung} der Eigenwerte und deren \emph{Vorzeichenstruktur}
in einer indefiniten Metrik (z.B. durch eine Pseudometrik mit
Signatur $(1,3)$, die aus der Gram-Matrix durch eine nicht-triviale
Bilinearform induziert wird). Aber das erfordert einen indefiniten
Kernel -- und der ist in PST-2-M nicht definiert.}
\begin{Fix}
In PST-3-P, Sektion 7, die Hypothese überarbeiten oder
streichen. Als Ersatz eine neue Hypothese einführen, die auf der
\textbf{konstruktiv/destruktiv-Idee} basiert (siehe Fixing-Vorschlag
zu \textbf{[K-15]}). Die alte Hypothese als ``verworfener Ansatz''
in eine Methodennote verlagern.
\end{Fix}

%--------------------------------------------------------------
\subsection{\K{08} PST-4-G: Die zentrale Behauptung $d_{\text{eff}} \approx 4$ ist nicht verifizierbar}
%--------------------------------------------------------------

\Dok{PST-4-G, Sektion 5}
\Stelle{Befund ``Zentraler PST-Befund: $d_{\text{eff}} \approx 4$''}
\Prob{Das Dokument behauptet:
\begin{quote}
``Diffusion Maps mit $t \approx 9$--$10$ liefern reproduzierbar:
$d_{\text{eff}} = 3.98 \pm 0.36$ über verschiedene OPP-Stichproben
($N = 1\,000$ bis $N = 100\,000$), verschiedene Kernel-Funktionen,
verschiedene $K$-Werte ($K = 10$ bis $K = 50$), verschiedene
Schwellenwerte $\theta$ ($0.90$ bis $0.99$).''
\end{quote}
Aber: \textbf{Es gibt keinen Code, keine Rohdaten, keine
GitHub-Links.} Die einzige verifizierbare Simulation, die wir bisher
gesehen haben, ist \textbf{Test 35} ($N = 40$, 2D-Punkte, 10 Runs),
bei dem $d_{\text{eff}}$ bei $t = 10$ exakt bei $3.98$ lag -- aber als
Teil eines \textbf{monotonen Abfalls} ($t = 6 \to 7.8$, $t = 10 \to
4.0$, $t = 12 \to 3.1$).}
\Krit{Ein wissenschaftliches Dokument, das reproduzierbare Ergebnisse
beansprucht, muss die \textbf{Reproduktionsgrundlage} liefern. Ohne
Code und Rohdaten ist das ein \textbf{Glaubensbekenntnis}, kein
Befund. Ein externer Reviewer (oder eine zukünftige KI) kann diese
Behauptung nicht prüfen.}
\Empf{Entweder (a) den Code und die Rohdaten sofort auf GitHub
veröffentlichen (mit Commit-Hash, Zeitstempel, verwendeter
Hardware) -- oder (b) die Behauptung auf ``Vorläufiger Befund aus
internen Tests, Dokumentation folgt'' abschwächen.}
\begin{Fix}
In PST-4-G, Sektion 5, einen neuen Absatz einfügen:
\begin{verbatim}
\begin{methodennote}[Reproduzierbarkeit und Datenlage]
Die hier berichteten Ergebnisse stammen aus internen
Testreihen, die derzeit nicht vollständig versioniert sind.
Eine vollständige Dokumentation inklusive Code, Rohdaten,
Parameter-Metadaten und Reproduktions-Skripten ist in
Vorbereitung (siehe Framework-Dokumentation, geplant für
PST-7-F). Bis zur Veröffentlichung dieser Daten gilt der
Befund als vorläufig.
\end{methodennote}
\end{verbatim}
\end{Fix}

%--------------------------------------------------------------
\subsection{\K{09} PST-4-G: Die Monotonie von $d_{\text{eff}}(t)$ wird nicht thematisiert}
%--------------------------------------------------------------

\Dok{PST-4-G, Sektion 5}
\Stelle{Remark ``Bedeutung von $t \approx 9$--$10$''}
\Prob{Das Dokument beschreibt:
\begin{itemize}
  \item Kleines $t$: $d_{\text{eff}}$ streut stark
  \item Mittleres $t$ (9--10): $d_{\text{eff}}$ konvergiert stabil auf $\approx 4$
  \item Großes $t$ (20--50): $d_{\text{eff}}$ nimmt wieder ab
\end{itemize}
Aber es \textbf{zeigt nicht die vollständige Kurve}. Wir wissen aus
Test 35, dass $d_{\text{eff}}$ monoton mit $t$ sinkt. Die
``Konvergenz bei $t \approx 9$--$10$'' könnte ein \textbf{Plateau}
in diesem Abfall sein -- oder ein \textbf{Artefakt} der spezifischen
Parameter ($N$, Kernel, $\theta$).}
\Krit{Wenn $d_{\text{eff}}(t)$ eine glatte, monotone Funktion ist,
dann ist der Wert 4 bei $t = 10$ nicht fundamental. Er ist ein
\textbf{Schnitt durch eine kontinuierliche Kurve}. Die Physik erfordert
aber eine \textbf{fundamentale} Dimension -- nicht eine, die von einem
freien Parameter $t$ abhängt.}
\Empf{Die vollständige Kurve $d_{\text{eff}}(t)$ für
verschiedene $N$ und Kernel dokumentieren. Wenn es ein Plateau gibt,
zeigen. Wenn es keines gibt, ehrlich sagen: ``Der Wert 4 ist
parameterabhängig und nicht fundamental.''}
\begin{Fix}
In PST-4-G, Sektion 5, die Tabelle der Test-35-Ergebnisse
vollständig einfügen (alle $t$-Werte von 6 bis 12 mit
Standardabweichungen). Hinzufügen: ``Die Monotonie von
$d_{\text{eff}}(t)$ zeigt, dass der Wert 4 kein fundamentales
Merkmal des OPP ist, sondern eine Eigenschaft des
Diffusion-Map-Operators bei bestimmten Parametern. Die Frage,
ob es ein physikalisch ausgezeichnetes $t$ gibt, bleibt offen.''
\end{Fix}

%--------------------------------------------------------------
\subsection{\K{10} PST-4-G: $N = 100\,000$ ist unglaubwürdig ohne Kontext}
%--------------------------------------------------------------

\Dok{PST-4-G, Sektion 5.2}
\Stelle{Befund ``Reproduzierbarkeit''}
\Prob{Die Behauptung $N = 100\,000$ wirft Fragen auf:
\begin{itemize}
  \item $N = 100\,000$ Zustände mit einer dichten Gram-Matrix
    erfordern \textbf{80 GB RAM} ($100\,000^2 \times 8$ Byte).
  \item Die Eigenwertberechnung einer $100\,000 \times 100\,000$
    Matrix ist \textbf{rechenintensiv} ($O(N^3)$, also $\sim 10^{15}$
    Operationen).
  \item Auf einem i7-10700K würde das \textbf{Tage} dauern.
  \item Der Autor hat gesagt, die Rechner sind noch nicht aufgebaut.
\end{itemize}}
\Krit{Die Behauptung $N = 100\,000$ ohne Kontext ist eine
\textbf{Glaubensbehauptung}. Wurde das auf Grok ausgeführt? (Grok
hat Zeitlimits von Minuten, nicht Tagen.) Wurde das auf Cloud-Hardware
ausgeführt? (Dann: Welche? Kosten?) Oder ist $N = 100\,000$ eine
\textbf{Projektion} für zukünftige Tests?}
\Empf{Explizit benennen, welche $N$ auf welcher Hardware mit welcher
Laufzeit getestet wurden. ``$N = 100\,000$'' ohne Kontext ist eine
Behauptung, die das Vertrauen untergräbt.}
\begin{Fix}
In PST-4-G, Sektion 5.2, eine Tabelle einfügen:
\begin{verbatim}
\begin{center}
\begin{tabular}{llll}
\toprule
\textbf{N} & \textbf{Hardware} & \textbf{Laufzeit} & \textbf{Status} \\
\midrule
40         & Grok (Cloud)      & $<$ 1 Min.        & Dokumentiert \\
1\,000     & [noch offen]      & [noch offen]      & Geplant \\
10\,000    & [noch offen]      & [noch offen]      & Geplant \\
100\,000   & i7-10700K + 256GB & geschätzt 3--5 Tage & Zukünftig \\
\bottomrule
\end{tabular}
\end{center}
\end{verbatim}
\end{Fix}

%--------------------------------------------------------------
\subsection{\K{11} PST-5-S: Der ``harte Kern'' ist zu hart}
%--------------------------------------------------------------

\Dok{PST-5-S, Sektion 1}
\Stelle{Definition ``Harter Kern der PST''}
\Prob{Der vierte Punkt des harten Kerns lautet: ``Es gibt einen
universellen Projektionsoperator $\mathcal{P}^*$.'' Das ist eine
\textbf{massive Annahme}. Wenn sich herausstellt, dass es keinen
\emph{einen} universellen Operator gibt, sondern eine \textbf{Klasse}
von kontextabhängigen Operatoren (je nach Beobachter, Energieskala,
etc.), würde das die PST nicht zwingend scheitern lassen. Es
würde nur bedeuten, dass $\mathcal{P}^*$ kein einzelner Operator,
sondern ein \textbf{Funktor} oder eine \textbf{Familie} ist.}
\Krit{In Lakatos' Terminologie gehört in den harten Kern nur das,
was bei Aufgabe das Programm beendet. Die Existenz eines \emph{einzigen}
universellen Operators ist stärker als nötig. Die Aussage
``Raumzeit emergiert durch Projektion'' ist hart. Die Aussage ``Es
gibt genau ein $\mathcal{P}^*$'' ist eine Hilfshypothese im
Schutzgürtel.}
\Empf{Punkt 3 des harten Kerns ändern zu: ``Es gibt eine Klasse von
Projektionsoperatoren, die aus dem OPP Raumzeit emergieren lassen.''
Und $\mathcal{P}^*$ als spezifische Hypothese in den Schutzgürtel
verschieben.}
\begin{Fix}
In PST-5-S, Sektion 1, die Definition des harten Kerns
ändern:
\begin{verbatim}
\begin{definition}[Harter Kern der PST -- korrigiert]
\begin{enumerate}
  \item Der OPP existiert als hochdimensionaler, prä-metrischer
    Zustandsraum.
  \item Raumzeit, Physik und Bestimmtheit emergieren durch Projektion.
  \item Es gibt eine Klasse von Projektionsoperatoren, die aus dem OPP
    Raumzeit emergieren lassen.
  \item Das Prinzip der minimalen Ontologie gilt.
\end{enumerate}
\end{definition}
\end{verbatim}
\end{Fix}

%--------------------------------------------------------------
\subsection{\K{12} PST-5-S und PST-6-R: Der Weg zur Lorentz-Signatur über
komplexe Gram-Matrizen ist eine Sackgasse}
%--------------------------------------------------------------

\Dok{PST-5-S, Sektion 4 (Priorität 2); PST-6-R, Sektion 2}
\Stelle{``Imaginäre Eigenwerte bei komplexen Gram-Matrizen'' und
``Wick-Rotation''}
\Prob{Dieser Ansatz taucht in drei Dokumenten auf (PST-3-P, PST-5-S,
PST-6-R) und ist mathematisch mit der in PST-2-M festgelegten
reell-symmetrischen Gram-Matrix \textbf{unvereinbar}.
\begin{itemize}
  \item Eine komplexwertige Gram-Matrix verliert die Garantie
    reeller Eigenwerte und positiver Semidefinitheit.
  \item Die Wick-Rotation in der QFT setzt \textbf{bereits} die
    Minkowski-Metrik voraus -- sie ist ein Trick \emph{innerhalb}
    der Physik, nicht eine Herleitung \emph{der} Physik.
  \item Der Ansatz verfestigt sich in der Forschungsagenda, obwohl
    er ein mathematischer Kategorienfehler ist.
\end{itemize}}
\Krit{Die PST verfolgt einen mathematisch unhaltbaren Ansatz und
ignoriert gleichzeitig eine vielversprechendere Alternative (siehe
\textbf{[B-01]}). Das ist eine \textbf{Fehlallokation von
Aufmerksamkeit} mit potenziell fatalen Konsequenzen für die
weitere Entwicklung.}
\Empf{Diesen Ansatz aus der Prioritätenliste streichen oder als
``hochspekulativ, erfordert Neufundierung der Gram-Matrix-Theorie''
kennzeichnen. Stattdessen den Weg über die
\textbf{konstruktiv/destruktive Dimension} (Raum = konstruktiv, Zeit
= destruktiv) als primären Kandidaten nennen.}
\begin{Fix}
In PST-5-S, Sektion 4, Priorität 2 ersetzen durch:
\begin{verbatim}
\begin{ziel}[Priorität 2: Lorentz-Signatur -- korrigiert]
\textbf{Frage:} Wie emergiert $(+,-,-,-)$ aus einer signaturfreien
Gram-Matrix?

\textbf{Kandidatenansatz (neu):}
Die drei Raumdimensionen sind ``konstruktiv'' -- sie stärken die
Bildung von Attraktoren. Die Zeitdimension ist ``destruktiv''
-- sie zerstört Attraktoren (Entropie, Informationsverlust).
Die Signatur $(+,-,-,-)$ reflektiert diesen ontologischen Gegensatz.
Mathematisch könnte dies durch eine indefinite Bilinearform
auf der Gram-Matrix realisiert werden, bei der konstruktive Eigenwerte
positiv und destruktive Eigenwerte negativ sind.

\textbf{Verworfener Ansatz:} Imaginäre Eigenwerte bei komplexen
Gram-Matrizen (inkonsistent mit PST-2-M; siehe Review [K-12]).
\end{ziel}
\end{verbatim}
\end{Fix}

%--------------------------------------------------------------
\subsection{\K{13} PST-6-R: Wick-Rotation ist physikalisch irreführend}
%--------------------------------------------------------------

\Dok{PST-6-R, Sektion 2}
\Stelle{Hypothese ``Wick-Rotation als projektive Eigenschaft''}
\Prob{Die Wick-Rotation $t \to -i\tau$ in der QFT ist ein
\textbf{mathematisches Hilfsmittel}, keine physikalische
Transformation. Sie macht das Pfadintegral konvergent, indem sie
von Minkowski- zu Euklidischer Zeit wechselt. Aber: Die QFT setzt
\textbf{bereits} die Minkowski-Metrik voraus.}
\Krit{Die PST behauptet, die Lorentz-Signatur aus dem OPP
herzuleiten. Wenn sie dafür die Wick-Rotation verwendet, setzt sie
implizit voraus, was sie beweisen will: die Existenz einer
Zeitkoordinate $t$, die man rotieren kann. Das ist ein
\textbf{petitio principii}.}
\Empf{Streichen oder als ``heuristische Analogie, keine Herleitung''
kennzeichnen. Die Wick-Rotation kann nicht das Fundament sein, auf dem
die Lorentz-Signatur emergiert.}
\begin{Fix}
In PST-6-R, Sektion 2, die Hypothese umbenennen in
\texttt{remark} mit dem Text: ``Die Wick-Rotation in der QFT ist ein
Standardtrick zur Konvergenz des Pfadintegrals. Sie setzt jedoch
bereits die Minkowski-Metrik voraus und kann daher nicht als
Herleitung der Lorentz-Signatur aus dem OPP dienen. Sie bleibt eine
heuristische Analogie.''
\end{Fix}

%--------------------------------------------------------------
\subsection{\K{14} PST-6-R: Die Schrödinger-Gleichung als
``projektive Eigenschaft'' ist zu vage}
%--------------------------------------------------------------

\Dok{PST-6-R, Sektion 4}
\Stelle{Hypothese ``Schrödinger-Gleichung aus Projektion''}
\Prob{Die Hypothese sagt: ``Die Schrödinger-Gleichung könnte
eine projektive Eigenschaft sein.'' Aber sie liefert keinen
Mechanismus. Die Schrödinger-Gleichung enthält:
\begin{itemize}
  \item Eine Ableitung nach der Zeit ($\partial/\partial t$)
  \item Einen Hamilton-Operator $\hat{H}$
  \item Die imaginäre Einheit $i$
\end{itemize}
Keines dieser Elemente ist in der PST bisher definiert.}
\Krit{Die Hypothese suggeriert Substanz, wo keine ist. Sie ist nicht
falsch -- aber sie ist leer. Das ist für ein Forschungsprogramm in
Ordnung, aber die Sprache sollte transparenter sein.}
\Empf{Entweder (a) explizit als ``noch völlig offen, keine
mathematische Substanz'' kennzeichnen, oder (b) einen konkreten
Ansatz skizzieren.}
\begin{Fix}
In PST-6-R, Sektion 4, die Hypothese durch eine
\texttt{methodennote} ersetzen:
\begin{verbatim}
\begin{methodennote}[Schrödinger-Gleichung -- offenes Terrain]
Die Herleitung der Schrödinger-Gleichung aus der Projektion
ist ein Fernziel der PST. Bisher existiert kein mathematischer
Mechanismus, der die Elemente der Gleichung ($i$, $\hat{H}$,
$\partial/\partial t$) aus PST-Strukturen ableitet. Dieser
Abschnitt kartografiert das Terrain, nicht den Weg.
\end{methodennote}
\end{verbatim}
\end{Fix}

\newpage

%==============================================================
\section{Bedeutsame Befunde [B]}
%==============================================================

%--------------------------------------------------------------
\subsection{\B{01} PST-1-F bis PST-6-R: Die destruktive Zeitdimension
fehlt komplett}
%--------------------------------------------------------------

\Dok{Alle PST-Dokumente}
\Stelle{Fehlend}
\Prob{Die stärkste originelle Idee des gesamten OP-Programms --
die konstruktiven Raumdimensionen vs. die destruktive Zeitdimension
als ontologische Begründung der Metrik-Signatur -- taucht in
keinem PST-Dokument auf.}
\Krit{Die PST verfolgt einen mathematischen Weg (komplexe
Eigenwerte, Wick-Rotation), der problematisch ist, während sie
einen philosophisch fundierten Weg (konstruktiv/destruktiv)
ignoriert, der mathematisch sauberer wäre. Das ist ein
\textbf{Blindfleck}.}
\Empf{Diese Idee muss in eine Überarbeitung der PST
einfließen. Sie könnte die Lorentz-Signatur ontologisch
begründen, ohne die mathematische Konsistenz der Gram-Matrix zu
zerstören.}
\begin{Fix}
In PST-5-S, Sektion 4, als neue Priorität einfügen:
\begin{verbatim}
\begin{ziel}[Priorität 2b: Ontologische Begründung der
Lorentz-Signatur]
\textbf{Hypothese:} Die drei Raumdimensionen sind ``konstruktiv''
-- sie stärken die Bildung und Stabilität von Attraktoren.
Die Zeitdimension ist ``destruktiv'' -- sie zerstört Attraktoren
(Entropie, Informationsverlust). Die Signatur $(+,-,-,-)$
reflektiert diesen ontologischen Gegensatz.

\textbf{Mathematische Umsetzung:} Eine indefinite Bilinearform
auf der Gram-Matrix, bei der konstruktive Eigenwerte
positiv und destruktive Eigenwerte negativ sind. Dies würde
eine indefinite Metrik mit Signatur $(1,3)$ oder $(3,1)$ erzeugen,
ohne die reelle Struktur der Gram-Matrix aufzugeben.
\end{ziel}
\end{verbatim}
In PST-6-R, Sektion 2, eine neue Hypothese einführen, die diesen
Ansatz entwickelt.
\end{Fix}

%--------------------------------------------------------------
\subsection{\B{02} PST-3-P: Die Verbindung $I_0 \leftrightarrow \hbar$
bleibt dimensionsanalytisch offen}
%--------------------------------------------------------------

\Dok{PST-3-P, Sektion 6}
\Stelle{Hypothese ``Quantisierung aus endlicher Projektionsauflösung''}
\Prob{$[I_0] = \text{Bit}$ (Informationsmaß), $[\hbar] =
\text{J} \cdot \text{s} = \text{kg} \cdot \text{m}^2/\text{s}$
(Wirkung). Die Hypothese behauptet eine Korrespondenz, aber ohne
eine \textbf{Konstante}, die Bits in Wirkung umrechnet, ist das eine
Typengleichung, keine physikalische Gleichung.}
\Krit{In der Physik ist jede Gleichung dimensionshomogen. Wenn
$I_0 = \hbar$ gelten soll, braucht man einen Faktor $k$ mit
$[k] = \text{J} \cdot \text{s}/\text{Bit}$. Woher kommt dieser
Faktor? Aus dem OPP? Aus der Projektion? Das ist die entscheidende
Frage.}
\Empf{Explizit als offene Frage formulieren: ``Die Korrespondenz
$I_0 \leftrightarrow \hbar$ erfordert eine dimensionsanalytische
Brücke, die derzeit nicht vorhanden ist.''}
\begin{Fix}
In PST-3-P, Sektion 6, die Hypothese ergänzen:
\begin{verbatim}
\begin{methodennote}[Dimensionsanalytische Lücke]
Die Korrespondenz $I_0 \leftrightarrow \hbar$ erfordert eine
Brückenkonstante $k$ mit $[k] = \text{J} \cdot
\text{s}/\text{Bit}$. Ein möglicher Ansatz: $I_0 = \hbar /
(k_B T \cdot \tau)$, wobei $T$ und $\tau$ emergente Größen
der Projektion wären. Dies bleibt ein offenes Problem.
\end{methodennote}
\end{verbatim}
\end{Fix}

%--------------------------------------------------------------
\subsection{\B{03} PST-3-P: Die ``Proposition'' Beobachter =
Projektionsoperatoren ist ein Postulat}
%--------------------------------------------------------------

\Dok{PST-3-P, Sektion 8}
\Stelle{Proposition ``Beobachter = Projektionsoperatoren''}
\Prob{Eine Proposition in der Mathematik ist eine Behauptung, die
aus Definitionen und Axiomen \emph{bewiesen} werden kann. Hier wird
behauptet: ``Ein physikalischer Beobachter ist identisch mit einem
Projektionsoperator.'' Das ist keine Folgerung -- das ist eine
\textbf{Setzung}, ein Postulat.}
\Krit{Die Proposition suggeriert Beweisbarkeit, wo keine ist. Das
ist ein Verstoß gegen die MKO-Sprachregel (Ebene 1 statt Ebene
3/5).}
\Empf{Umbenennen in \texttt{postulat} oder \texttt{definition}. Oder:
Als \texttt{proposition} belassen, aber mit einem Beweis versehen,
der aus den Axiomen des OP folgt.}
\begin{Fix}
In PST-3-P, Sektion 8, die \texttt{proposition} durch
\texttt{postulat} ersetzen und hinzufügen: ``Dies ist eine
grundlegende Setzung der PST, keine ableitbare Proposition. Sie
verbindet die philosophische Kategorie des Beobachters (DPR) mit
der mathematischen Kategorie des Projektionsoperators.''
\end{Fix}

%--------------------------------------------------------------
\subsection{\B{04} PST-4-G: Die Traum-Analogie / SSD-Analogie fehlt}
%--------------------------------------------------------------

\Dok{PST-4-G, Gesamtdokument}
\Stelle{Fehlend}
\Prob{Das Dokument beschreibt den Paradigmenwechsel von ``OPP hat
Dimensionen'' zu ``Projektion erzeugt Dimensionen''. Aber es
vermisst die \textbf{intuitive Verankerung}, die der Autor dem
Reviewer erklärt hat: Die SSD-Analogie (Bits $\to$ Player $\to$
Film).}
\Krit{Die SSD-Analogie ist die beste didaktische Brücke zwischen
Philosophie und Mathematik. Ihr Fehlen macht PST-4-G für externe
Leser schwerer zugänglich.}
\Empf{Einen Abschnitt ``Intuitive Verankerung'' hinzufügen.}
\begin{Fix}
In PST-4-G, nach Sektion 2 (Phase 2), eine neue Sektion
einfügen:
\begin{verbatim}
\begin{remark}[Intuitive Verankerung: Die SSD-Analogie]
Der OPP ist wie eine SSD: statisch, zeitlos, nur Bits.
Der Diffusion-Map-Operator ist wie ein Player: Er projiziert
sequentiell und erzeugt emergent Raum und Zeit.
$d_{\text{eff}} \approx 4$ ist die emergente Dimension des Films.
Diese Analogie ist keine mathematische Herleitung, sondern ein
didaktisches Hilfsmittel (MKO Ebene 4).
\end{remark}
\end{verbatim}
\end{Fix}

%--------------------------------------------------------------
\subsection{\B{05} PST-5-S: Die Verbindung zu Connes ist konträr,
nicht komplementär}
%--------------------------------------------------------------

\Dok{PST-5-S, Sektion 7}
\Stelle{``PST und nicht-kommutative Geometrie (Connes)''}
\Prob{Connes' spektrale Tripel $(A, H, D)$ setzt einen
\textbf{Dirac-Operator $D$} voraus, der eine \textbf{Metrik} kodiert.
Die PST behauptet, Metrik emergiere erst durch Projektion. Das sind
\textbf{konträre Ausgangspunkte}, nicht komplementäre.}
\Krit{Wenn die PST sagt ``Metrik ist emergent'', dann kann sie
nicht gleichzeitig Connes' Ansatz als ``eine mathematisch
ausgearbeitete Version des PST-Projektionsprinzips'' interpretieren.
Connes \emph{setzt} die Metrik durch $D$ voraus; die PST
\emph{leitet} sie ab.}
\Empf{Die Verbindung zu Connes als ``interessante mathematische
Analogie, aber mit umgekehrter ontologischer Richtung''
kennzeichnen -- oder streichen.}
\begin{Fix}
In PST-5-S, Sektion 7, den Text ändern:
\begin{verbatim}
\begin{remark}[PST und nicht-kommutative Geometrie (Connes)]
Alain Connes beschreibt Raumzeit über spektrale Tripel
$(A, H, D)$. Der Dirac-Operator $D$ kodiert dabei bereits eine
Metrik. Die PST und Connes' Ansatz sind daher ontologisch
konträr: Connes setzt Metrik voraus, die PST leitet sie ab.
Ob beide Programme letztlich formal äquivalent sind, ist eine
offene Frage. Die PST kann Connes' Mathematik als Werkzeug
nutzen, nicht als ontologische Grundlage.
\end{remark}
\end{verbatim}
\end{Fix}

%--------------------------------------------------------------
\subsection{\B{06} PST-5-S: Das Scheitern-Kriterium 3 ist mathematisch trivial}
%--------------------------------------------------------------

\Dok{PST-5-S, Sektion 5}
\Stelle{``Wann scheitert die PST?'' -- Kriterium 3}
\Prob{``Die Lorentz-Signatur grundsätzlich nicht aus einer reellen,
positiv semidefiniten Gram-Matrix emergieren kann.'' Das ist kein
mögliches Scheitern -- das ist ein \textbf{mathematisches
Faktum}. Eine positiv semidefinite Matrix hat nur nicht-negative
Eigenwerte. Eine indefinite Metrik mit Signatur $(+,-,-,-)$ hat
positive und negative Eigenwerte. Das eine kann nicht aus dem
anderen folgen.}
\Krit{Das Scheitern-Kriterium suggeriert, dass dies eine offene
empirische Frage ist. Aber es ist eine geschlossene mathematische.
Die PST muss entweder (a) indefinite Gram-Matrizen zulassen, (b)
die Signatur durch einen anderen Mechanismus erzeugen, oder (c)
zugeben, dass die reell-positive Definitheit eine zu starke Annahme
war.}
\Empf{Das Kriterium umformulieren: ``Es zeigt sich, dass keine
Modifikation der Gram-Matrix-Struktur (z.B. Übergang zu indefiniten
Bilinearformen) die Lorentz-Signatur konsistent erzeugen kann.''}
\begin{Fix}
In PST-5-S, Sektion 5, Kriterium 3 ersetzen durch:
\begin{verbatim}
\item Es zeigt sich, dass keine konsistente Modifikation der
  Gram-Matrix-Struktur (einschließlich indefiniter
  Bilinearformen) die Lorentz-Signatur $(+,-,-,-)$ erzeugen
  kann.
\end{verbatim}
\end{Fix}

%--------------------------------------------------------------
\subsection{\B{07} PST-6-R: Die Verbindung zur ART ist zu oberflächlich}
%--------------------------------------------------------------

\Dok{PST-6-R, Sektion 6}
\Stelle{``Allgemeine Relativitätstheorie (ART)''}
\Prob{Die Behauptung, die Einstein-Feldgleichungen ``könnten aus
Eigenschaften des Projektionsspektrums emergieren'', ist nicht
falsch -- aber sie ist \textbf{leer}. Es wird nicht erklärt:
\begin{itemize}
  \item Wie $G_{\mu\nu}$ (Einstein-Tensor) aus einer Gram-Matrix folgen
    soll
  \item Wie $T_{\mu\nu}$ (Energie-Impuls-Tensor) als
    Projektionseigenschaft interpretiert werden soll
  \item Wie die Kovarianz (allgemeine Koordinatentransformationen)
    emergieren soll
\end{itemize}}
\Krit{Die Sprache suggeriert mehr Substanz als vorhanden. Das ist
im Rahmen eines Forschungsprogramms legitim -- aber die Sprache
sollte transparenter sein zwischen ``intuitiver Analogie'' und
``mathematischer Hypothese''.}
\Empf{Entweder einen konkreten mathematischen Ansatz skizzieren
oder als ``völlig offen'' kennzeichnen.}
\begin{Fix}
In PST-6-R, Sektion 6, den Text ändern:
\begin{verbatim}
\begin{remark}[Allgemeine Relativitätstheorie (ART) -- offen]
Die Verbindung zwischen PST und ART ist völlig offen.
Insbesondere fehlt jeder mathematische Mechanismus, der:
\begin{itemize}
  \item den Einstein-Tensor $G_{\mu\nu}$ aus einer Gram-Matrix
    ableitet,
  \item den Energie-Impuls-Tensor $T_{\mu\nu}$ als
    Projektionseigenschaft interpretiert,
  \item die allgemeine Kovarianz aus der Projektionsstruktur
    emergieren lässt.
\end{itemize}
Dieser Abschnitt kartografiert das Terrain, nicht den Weg.
\end{remark}
\end{verbatim}
\end{Fix}

\newpage

%==============================================================

\vspace{2em}
\noindent\rule{\textwidth}{0.4pt}
\vspace{1em}

\noindent\textbf{Hinweis:} Die \textbf{minoren Befunde [M]},
\textbf{positiven Bewertungen [P]} und die \textbf{priorisierten
Empfehlungen} für die Überarbeitung sind in \textbf{Teil~2}
dieses Reviews dokumentiert.

\vspace{2em}

\begin{center}
\rule{0.5\textwidth}{0.4pt}
\end{center}

\vspace{1em}

\noindent\textit{Dieses Review wurde gemäß dem Adversarial Review
Protocol (AR) verfasst. Der Reviewer hat die Probleme benannt und
Lösungsvorschläge skizziert, aber die inhaltliche Entscheidung
und Umsetzung obliegt dem Autorenteam.}

\end{document}
